\documentclass[11pt,a4paper]{article}
\usepackage[utf8]{inputenc}
\usepackage[T1]{fontenc}
\usepackage[a4paper,margin=2.5cm]{geometry}
\usepackage{microtype}
\usepackage{hyperref}
\usepackage{xurl}
\usepackage{xcolor}
\sloppy
\hypersetup{colorlinks=true, linkcolor=blue, urlcolor=blue, citecolor=blue}

\setlength{\parskip}{0.4em}
\setlength{\parindent}{0pt}

\begin{document}

\textbf{Cover letter to the Editors of \emph{Physical Review A}}

\vspace{0.4cm}

Hikaru Wakaura, Taiki Tanimae \\
QIRI (Quantum Integrated Research Institute Inc.)\\
1--16--3, Akasaka, Minato-ku, Tokyo 107--0061, Japan\\
\href{mailto:h.wakaura@qiri.co.jp}{h.wakaura@qiri.co.jp} \quad ORCID: 0009-0001-2341-7710

\vspace{0.3cm}
October 9, 2026

\vspace{0.4cm}

Dear Editors,

\medskip

We are pleased to submit our manuscript, \textbf{``Catalytic quantum error correction: efficient catalyst preparation and a completely positive covariant recovery channel''}, for consideration as a Regular Article in \emph{Physical Review A} (Quantum information).

\medskip

\textbf{Summary of the contribution.}
Shiraishi and Takagi (\emph{Phys.\ Rev.\ Lett.}\ \textbf{132}, 180202, 2024) proved that catalytic covariant operations can amplify quantum coherence at an unbounded asymptotic rate.
Turning that theorem into an operational state-recovery scheme has been blocked by two obstacles: the constructive proof requires $n^{*} \sim d^{4} e^{2\gamma}$ noisy copies to prepare the catalyst (about $7.3 \times 10^{10}$ for six qubits at dephasing $\gamma = 2$), and no concrete covariant recovery channel had been exhibited.
This manuscript addresses both.
(i)~A three-stage pipeline---CPMG dynamical decoupling, Clifford twirling, recursive swap-test purification---prepares catalysts at the matched $F_\mathrm{cat} \geq 0.99$ endpoint from 32~copies, a $10^{4}$--$10^{9}$-fold reduction depending on dimension, with postselection costs quantified ($P_\mathrm{succ} \approx 0.7$).
(ii)~An explicitly completely positive, exactly covariant joint-channel implementation of the recovery map (system--catalyst--ancilla; catalyst obtained by partial trace, never copied) achieves genuine fidelity recovery under dephasing at $d = 4$ ($0.54 \to 0.77$), passes a target-swap control experiment, and delimits the method's current reach: shallow covariant circuits do not yet help under depolarizing noise or at $d = 8$, and they consume rather than preserve the catalyst.
An effective surrogate model of the asymptotic limit, benchmarked across 200~noise configurations ($d = 4$--$64$), is clearly separated from the channel-level results and collected in an appendix.

\medskip

\textbf{Why \emph{Physical Review A}.}
\begin{itemize}
\item \emph{Protocol with explicit success conditions.} The CQEC protocol is formalized through the mode-inclusion criterion $\mathcal{C}(\rho_{0}) \subseteq \mathcal{C}(\rho_\mathrm{noisy})$, with finite-copy fidelity bounds and a clear statement of the operational assumptions (known target state, multiple copies), including its relation to stabilizer QEC, covariant codes, and the concurrent Purification QEC of Raghoonanan \& Byrnes (arXiv:2603.11568).
\item \emph{Verified numerics.} Every quantitative property of the joint channel is locked in by unit tests: $\|[U, H_\mathrm{total}]\| = 0$ exactly, linearity to machine precision, Choi positivity, trace preservation, and catalyst-by-partial-trace.
\item \emph{Scoped claims.} Model-level and channel-level results are explicitly distinguished throughout; sensitivity of the pipeline's dimension independence to the twirl model is derived rather than hidden.
\item \emph{Open science.} The reference implementation is released as the \texttt{cqec} Python package v0.2.0 (MIT license) with a 52-test \texttt{pytest} suite and a one-command reproduction script: \href{https://github.com/deeptell-inc/cqec}{\texttt{github.com/deeptell-inc/cqec}}.
\end{itemize}

\medskip

\textbf{Suggested referees} (institutional emails from public faculty pages; please verify before contact):
\begin{itemize}
\item \textbf{Naoto Shiraishi} (The University of Tokyo; \href{mailto:shiraishi@phys.c.u-tokyo.ac.jp}{\texttt{shiraishi@phys.c.u-tokyo.ac.jp}}) --- co-author of the foundational PRL on catalytic coherence transformations.
\item \textbf{Ryuji Takagi} (Nanyang Technological University; \href{mailto:ryuji.takagi@ntu.edu.sg}{\texttt{ryuji.takagi@ntu.edu.sg}}) --- co-author of the same PRL; resource-theoretic content.
\item \textbf{Tim Byrnes} (NYU Shanghai; \href{mailto:tim.byrnes@nyu.edu}{\texttt{tim.byrnes@nyu.edu}}) --- co-author of the concurrent Purification QEC framework compared in our taxonomy table.
\item \textbf{Iman Marvian} (Duke University; \href{mailto:iman.marvian@duke.edu}{\texttt{iman.marvian@duke.edu}}) --- covariant operations and resource theory of asymmetry.
\item \textbf{Andrew M.~Childs} (University of Maryland; \href{mailto:amchilds@umd.edu}{\texttt{amchilds@umd.edu}}) --- recursive swap-test purification adapted in our pipeline.
\end{itemize}

\medskip

\textbf{Statements.}
This work is not under consideration elsewhere.
A complementary manuscript, ``Blind Catalytic Quantum Error Correction'' (independent content, no shared figures or theorems), is in review at a different journal.
All authors have approved the submission; we declare no competing interests.
The manuscript is archived at \href{https://doi.org/10.48550/arXiv.2603.25774}{arXiv:2603.25774}.

\medskip

We thank you for your consideration and look forward to the referees' feedback.

\medskip

Sincerely,

\vspace{0.4cm}
Hikaru Wakaura\\
Taiki Tanimae\\
QIRI Inc., Tokyo

\end{document}
